% ---
% Reemplaza líneas 1139-1190 de Docs/Technical_Report.tex
% Decisiones: D-014, D-040, H-019, H-052
% ---

\subsubsection{Medición y valores de expectativa}

La medición en computación cuántica extrae información clásica del estado cuántico
al precio de colapsar la superposición. Para un observable $\hat{O}$ representado
por un operador hermítico, el valor de expectativa cuantifica el promedio estadístico
obtenido al medir dicho observable sobre un ensamble de estados idénticos $|\psi\rangle$:
\begin{equation}
  \langle \hat{O} \rangle = \langle \psi | \hat{O} | \psi \rangle.
  \label{eq:expectation}
\end{equation}
En el presente trabajo, los observables corresponden a las matrices
de Pauli $Z$ aplicadas a cada qubit individual $i \in \{1, \dots, n\}$. Su valor
de expectativa está en el intervalo $[-1, 1]$:
\begin{equation}
  \langle Z_i \rangle = \langle \psi | Z_i | \psi \rangle \in [-1, 1],
  \label{eq:expZ}
\end{equation}
donde $\langle Z_i \rangle = +1$ denota el estado computacional $|0\rangle$,
$\langle Z_i \rangle = -1$ denota el estado $|1\rangle$, y valores intermedios
describen superposiciones coherentes o mezclas estadísticas. El vector ordenado de
expectativas:
\begin{equation}
  \mathbf{e}_Z = \bigl(\langle Z_1 \rangle, \langle Z_2 \rangle, \dots, \langle Z_n \rangle\bigr)^\top \in [-1, 1]^n
  \label{eq:expZ_vector}
\end{equation}
resume la proyección del estado cuántico sobre cada qubit.

En un dispositivo físico cuántico, $\langle Z_i \rangle$ no se mide de forma
directa sino que se estima mediante el promedio de $n_{\text{shots}}$ mediciones
proyectivas independientes. Dado que los autovalores de $Z_i$ son $\pm 1$, cada disparo da $\pm 1$ con probabilidad $P(+1) = (1+\langle Z_i\rangle)/2$, de modo que la varianza de un disparo individual es $1-\langle Z_i\rangle^2$. La varianza del estimador empírico
$\hat{Z}_i = \frac{1}{n_{\text{shots}}} \sum_{s=1}^{n_{\text{shots}}} z_i^{(s)}$
está determinada exactamente por:
\begin{equation}
  \operatorname{Var}(\hat{Z}_i) = \frac{1 - \langle Z_i \rangle^2}{n_{\text{shots}}}.
  \label{eq:shot_noise_variance}
\end{equation}
La dispersión muestral es máxima cuando el estado es puramente ecuatorial
($\langle Z_i \rangle = 0$, varianza $1/n_{\text{shots}}$) y se anula si el qubit
se halla en un autoestado de la base computacional ($\langle Z_i \rangle = \pm 1$).

El núcleo de fidelidad cuántica, $K_Q(\mathbf{x}, \mathbf{x}') = |\langle \phi(\mathbf{x}) | \phi(\mathbf{x}') \rangle|^2$,
se estima en hardware de otra forma: se implementa el circuito de
\textit{compute-uncompute} $U^\dagger(\mathbf{x}') U(\mathbf{x})$ sobre el estado
$|0\rangle^{\otimes n}$ y se cuenta la frecuencia con la que se registra
la cadena nula $|0\dots 0\rangle$. Por lo tanto, la estimación del kernel en hardware
es un proceso binomial con probabilidad de éxito $p = K_Q$, cuya
varianza muestral es $K_Q(1 - K_Q)/n_{\text{shots}}$.

En la metodología principal del presente trabajo, con objeto de evaluar las
propiedades geométricas de las representaciones sin el ruido
que introduce un número finito de disparos, los valores de
expectativa y las matrices de fidelidad se evalúan analíticamente en simulación a
partir del vector de estado exacto (\textit{statevector}). El impacto empírico y
analítico de la cota finita de disparos ($n_{\text{shots}}$) se analiza de forma
explícita en el módulo de restricciones del simulador (M8).
